problem:  
Let \((M^{n+1},g)\) be a globally hyperbolic Lorentzian manifold satisfying the **null energy condition** and let \(\Sigma\subset M\) be a smooth, spacelike Cauchy hypersurface with induced Riemannian metric \(h\).  
Suppose \(\Sigma\) contains two disjoint, compact, two‑sided, **stable marginally outer trapped surfaces** \(S_1,S_2\subset\Sigma\), with \(S_2\) outside \(S_1\) in the outer normal direction.  
Assume the mean curvature of \(\Sigma\) is nonnegative and that each null congruence orthogonal to \(S_1\) has smooth optical scalars up to its first conjugate point.

For null geodesics \(\gamma:[0,L]\to M\) orthogonal to \(S_1\), denote by \(\theta_+(\lambda)\) their expansion satisfying the Raychaudhuri equation  
\[
\frac{d\theta_+}{d\lambda}=-\frac{1}{n-1}\theta_+^2-\sigma_+^2-\mathrm{Ric}(\ell_+,\ell_+),\qquad \theta_+(0)=0,
\]
and define the averaged **null Ricci integral** up to parameter length \(L>0\),
\[
\mathcal{R}(L)=\frac{1}{L}\int_0^L\mathrm{Ric}(\ell_+,\ell_+)\,d\lambda.
\]
Let \(d_\Sigma(S_1,S_2)\) be the Riemannian distance in \((\Sigma,h)\) between \(S_1\) and \(S_2\).

**Question:**  
Does there exist a universal constant \(C_n>0\) (depending only on dimension) such that whenever
\[
\mathcal{R}(L) > \frac{C_n}{\,d_\Sigma(S_1,S_2)^2\,},
\]
every null geodesic orthogonal to \(S_1\) develops a conjugate point before intersecting \(S_2\)?  
If so, can equality occur only in the shear‑free, Ricci‑flat case, implying that the segment of spacetime between \(S_1\) and \(S_2\) is locally isometric to a warped product
\[
([0,L]\times S_1,\; -2\,du\,dt + f(u)^2h_{S_1})?
\]

Equivalently: is there a **Lorentzian Raychaudhuri–Myers comparison principle** asserting that large *average null curvature* enforces a bound on the possible “spatial separation” of stable MOTSs, with equality forcing null rigidity?

---

Why is it a "good" problem:  
1. **Profound insight:**  
This problem conjectures a precise quantitative bridge between *null curvature focusing* and *spatial geometry* of marginally trapped surfaces.  It asks whether there exists a Lorentzian analogue of the classical Myers comparison theorem, but formulated in terms of null expansions and MOTS separation.  A positive answer would give a quasi‑local curvature–distance law for horizons, sharpening the geometric content of the Raychaudhuri equation.

2. **Bridge between fields:**  
It integrates **Lorentzian causality theory**, **geometric analysis of MOTSs**, and **comparison geometry**.  Addressing it would require reframing focusing inequalities in a Riemannian comparison spirit, linking local analytic behaviour of null congruences to global spacelike geometry—connecting general relativity, PDE analysis, and Riemannian-style curvature estimates.

3. **Extreme simplicity:**  
The statement is concise: two marginally trapped surfaces, their separation, and an inequality involving the averaged null Ricci curvature.  Despite its deceptively minimal form, the problem hides deep challenges concerning null geodesic focusing, stability operators, and the transition between local and global geometric regimes.

4. **Potential to open new directions:**  
Proving such an inequality would introduce genuine *null comparison geometry*, potentially yielding Lorentzian analogues of Myers’ diameter theorem or Bishop–Gromov bounds.  Failure of the inequality might uncover new counterexamples guiding our understanding of multiple‑horizon configurations and dynamical trapping phenomena.

This problem refines the conjecture to a mathematically coherent, testable inequality, at once simple and profound, creating a clear target for developing Lorentzian comparison geometry governed by marginally trapped surfaces.